\subsection{Working schemes via generated AECs and lifted types}
\label{subsec:alternate-working}
We now use Boney's alternate presentation as our definition. The source's
four-clause package is restated in Appendix~\ref{app:working-equivalence},
where we prove the comparison, with its precise scope restrictions.
This is a reorganization of the package, not a new verification of any
of the manuscript's applications.

\begin{definition}[Concrete coding convention]\label{def:working-coding}
Fix a nonempty finitely generated shape $\mathcal I$, an AEC $\K$, and
an unbounded definable class $\mathcal S$ of infinite cardinals with minimum
$\mu_0\geq\LS(\K)$. Use a fixed finitary many-sorted language
$\tau^{\mathrm{gen}}$, of size at most $\mu_0$, coding the vertex
structures, specified internal maps, and all fixed named data.
Its universe is the disjoint union of the vertex occurrences; no indexing
skeleton is added to this universe. Structural embeddings and isomorphisms
always refer to this entire code, including the names. Write $\|X\|$ for
its total cardinality and $|X|$ for its sorted universe.

An AEC $\K^{\mathrm{gen}}$ \emph{over}
$\operatorname{Diag}_{\mathcal I}(\K_{\geq\mu_0})$ means an AEC of these
codes, with no models below $\mu_0$, together with the concrete faithful
forgetful functor $U$ which drops the named data. Its strong embeddings
have underlying structural embeddings with strong vertex maps; the
converse is not required.
Sorted universes and maps are retained by $U$, and directed colimits are
computed on them, so $U$ preserves directed colimits. We do not require
$U$ to be full or injective on objects.

In this diagram version, the internal maps are structural embeddings.
For the source's more general specified maps, replace the target by the
category of finite systems of $\K$-models with arbitrary specified
functions and componentwise strong commuting embeddings. The same
sorted-code convention and the arguments below apply to that version.
This change of ambient category is necessary for a claim covering that
full generality of the source.
\end{definition}

\begin{definition}[Structural data and lifted automorphism comparison]
\label{def:working-scheme-data}
In the convention above, the data are:
\begin{enumerate}[label=\textup{(D\arabic*)},leftmargin=*]
\item An AEC $\K^{\mathrm{gen}}$ over the diagram category with
$\LS(\K^{\mathrm{gen}})=\mu_0$.
\item A set-sized finite-support term recipe $\Psi$ with values denoted
$F(L)=\operatorname{EM}_{\tau^{\mathrm{gen}}}(L,\Psi)$.
Here this notation denotes the \emph{system code}, not a structure
with an additional named skeleton. Elements have sorted labels on finite
increasing supports; equality and all symbols of the code are evaluated
uniformly from labels and finite order/equality patterns. This includes
fixed named data and compatibility with substitution of positions.
We require $F(\alpha)\in\K^{\mathrm{gen}}$ with every vertex of size
$\mu$ whenever $\mu\leq\alpha<\mu^+$ and $\mu\in\mathcal S$.
No membership assertion for arbitrary linear orders is included.
\item Let $\mathcal C_\mu$ be the structured-isomorphism closure of
these $F(\alpha)$, and put $\mathcal C=\bigcup_{\mu\in\mathcal S}\mathcal C_\mu$.
Specify
\[
 H(X,Y)\subseteq\operatorname{Emb}_{\K^{\mathrm{gen}}}(X,Y)
 \qquad(X,Y\in\mathcal C).
\]
Write $H_\mu$ for the same-level arrows, and require equality of both
objects and arrows
\[
 (\mathcal C_{\mu_0},H_{\mu_0})=\K^{\mathrm{gen}}_{\mu_0}.
\]
\end{enumerate}
For tuples of one fixed sorted index pattern, bounded in both containing
objects, define the relation $E$ as follows: $E((\ba;X),(\bb;Y))$ holds
if there exist an eligible $Z$, $f\in H(X,Z)$, $g\in H(Y,Z)$, and
$h\in\operatorname{Aut}(Z)$ such that $hf(\ba)=g(\bb)$.
Once $E$ is an equivalence relation, its classes are denoted
$\operatorname{tp}_{\mathrm{la}}(\ba;X)$.
All witnesses here are \emph{eligible} and all maps belong to $H$;
arbitrary $\K^{\mathrm{gen}}$-amalgams are not being used.
\end{definition}

\begin{definition}[Working scheme: alternate axioms]\label{def:working-scheme}
The preceding data form a working scheme if the following hold.
\begin{enumerate}[label=\textup{(A\arabic*)},leftmargin=*]
\item\label{ax:la-equivalence} \textbf{Equivalence.} For every bounded sorted index
pattern, $E$ is an equivalence relation across the eligible objects large
enough for that pattern. In particular, transitivity is a requirement,
not an automatic property of common-target comparisons.
\item\label{ax:order-lifts} \textbf{Presentation arrows.} Every order embedding
between well-orders of eligible, possibly different, cardinalities
induces an arrow in $H$.
\item\label{ax:structural-closure} \textbf{Structural and fixed-level closure.}
$H$ contains all structured isomorphisms and is closed under composition.
If $X\subseteq Y\subseteq Z$ are eligible and both inclusions into $Z$
belong to $H$, then $X\hookrightarrow Y$ belongs to $H$.
For each $\mu$, $(\mathcal C_\mu,H_\mu)$ is an AEC in $\mu$:
the induced strong-substructure relation is coherent, invariant under
isomorphism, and satisfies chain unions and smoothness at that size.
Explicitly, for a nonempty increasing chain of inclusion arrows of length
$0<\delta<\mu^+$, its union $X$ lies in $\mathcal C_\mu$, and every stage
includes into $X$ by an arrow. If the stages include by arrows into
\emph{any eligible} $Y$, their union includes into $Y$ by an arrow as well.
The larger-target clause is part of the axiom, not just same-level smoothness.
\item\label{ax:orbit} \textbf{Same-level orbit characterization.} For
$X,Y\in\mathcal C_\mu$ and tuples of the same sorted pattern of length
strictly below $\mu$,
\[
 \operatorname{tp}_{\mathrm{la}}(\ba;X)=
 \operatorname{tp}_{\mathrm{la}}(\bb;Y)
 \quad\Longleftrightarrow\quad
 \exists h:X\cong Y\ (h(\ba)=\bb).
\]
The isomorphism is of the whole structured system. Include the empty tuple.
\item\label{ax:exhaustion} \textbf{Countable exhaustion.} Let
$X\in\mathcal C_\mu$, $Y\in\mathcal C_\nu$, and let $p:X\to Y$ be a
structural embedding of codes, not initially assumed strong or in $H$.
Suppose there are increasing subsets $A_n\subseteq|X|$, $n<\omega$,
of size $<\mu$, with union $|X|$. For every vertex $v$, suppose also that
\[
 B^v_n\leq_\K B^v_{n+1}\leq_\K X(v),\qquad
 |B^v_n|<\mu,\qquad \bigcup_{n<\omega}B^v_n=X(v),\qquad
 |B^v_n|\subseteq A_n.
\]
Choose a sorted enumeration $\ba_n$ of $A_n$ and use its image enumeration
$p(\ba_n)$, with coordinate identifications consistent under inclusion.
If, for every $n$ and $v$,
\[
 \operatorname{tp}_{\mathrm{la}}(\ba_n;X)=
 \operatorname{tp}_{\mathrm{la}}(p(\ba_n);Y),\qquad
 p[B^v_n]\leq_\K Y(v),
\]
then $p\in H(X,Y)$. The sets $A_n$ need not be substructures.
The $B^v_n$ are additional witnesses contained in those sets; they need
not equal $A_n\cap X(v)$. Any extra geometric arrow condition must be
verified when establishing this axiom for a particular scheme.
\end{enumerate}
For comparison with the main regime of source Convention~5.1, take
eligible levels to be strong limits of countable cofinality. We keep
the countable exhaustion formulation explicit. At preliminary levels
of other cofinalities, equivalence with a more broadly indexed exhaustion
axiom is not asserted.
\end{definition}

\begin{remark}[Which clauses of the alternate draft became consequences?]
The draft separately lists mixed-arrow factorization, categoricity,
term naturality, and realization over actual tuples. We retain all four
as proved consequences below and in the appendix, rather than repeat
them as independent axioms. Their deletion from the axiom list does not
weaken the package. This also strengthens the draft's same-chart
naturality statement to the source's cross-level formulation.
\end{remark}

\begin{proposition}[Derived package]\label{prop:alternate-derived}
Under the alternate axioms:
\begin{enumerate}[label=\textup{(\roman*)},leftmargin=*]
\item There are common eligible targets. Lifted types include atomic
information, respect restrictions and reindexing, and are preserved by
every $H$-arrow. There is one empty type, so $\mathcal C_\mu$ is categorical
as a class of \emph{systems} at each eligible $\mu$.
\item Mixed downward closure holds: if $\nu\leq\mu$ are eligible,
every set of $\leq\nu$ sorted occurrences in $X\in\mathcal C_\mu$ lies
in an $H$-subobject in $\mathcal C_\nu$, retaining all fixed names.
\item For each fixed sorted arity $\theta$, the number of lifted types
has a set bound independent of level. A possible joint type is realized
over an actual tuple at a receiving level whenever the joint length is
bounded in both the witness and receiver.
\item Term tuples with the same labels and the same \emph{joint}
order/equality pattern of supports have equal lifted types, including
across eligible levels.
\end{enumerate}
\end{proposition}
\begin{proof}
See Lemmas~\ref{lem:la-basic}--\ref{lem:la-realization} in the appendix.
\end{proof}

\begin{proposition}[Mixed arrows from the family of level arrows]
\label{prop:mixed-factorization}
Let $X\in\mathcal C_\mu$, $Y\in\mathcal C_\nu$, $\mu<\nu$.
Every $f\in H(X,Y)$ factors as $j f_0$, where
$f_0\in H_\mu(X,Y_0)$ and $j:Y_0\to Y$ is, up to chart
isomorphisms, induced by inclusion of a suborder of size $\mu$ in a
chart for $Y$. Conversely, every such composite is in $H(X,Y)$.
Thus the \emph{family} $(H_\mu)_{\mu\in\mathcal S}$ and the presentation
arrows determine all mixed arrows of a working scheme.
\end{proposition}
\begin{proof}
Choose a chart for $Y$ and supports for all occurrences in $f[X]$.
Their union has size at most $\mu$; enlarge it to a suborder $J$ of size
$\mu$. Its presentation has image $Y_0$ containing $f[X]$, and its
inclusion $j$ is an arrow. The inclusion $f[X]\hookrightarrow Y$
is an arrow by isomorphism invariance. Coherence gives
$f[X]\hookrightarrow Y_0$ in $H$, hence $f_0\in H_\mu$.
The converse follows by composition. Fixed names are preserved by the
presentation arrows. This argument proves the factorization inside an
existing package; declaring this formula for arbitrary level data would
still require checking coherence and composition.
\end{proof}

\begin{remark}[The small-level completion and its remaining fullness question]
\label{rem:background-aec-generated}
The standard upward-completion theorem for an AEC in $\mu_0$ supplies
\[
 \K^{\mathrm{gen}}=(\mathcal C_{\mu_0},H_{\mu_0})^{\uparrow},
 \qquad \LS(\K^{\mathrm{gen}})=\mu_0.
\]
It recovers \emph{both objects and arrows} at $\mu_0$; see
Appendix~\ref{app:completion}. Earlier versions of these notes incorrectly
left all the AEC axioms of this completion as outstanding. The actual
unresolved issue here is higher-level fullness:
\[
 H(X,Y)\stackrel{?}{=}
 \operatorname{Emb}_{\K^{\mathrm{gen}}}(X,Y)
 \quad\text{for eligible }X,Y.
\]
Every designated arrow has a coherent small-arrow decomposition, but so
does every generated strong embedding. That observation alone does not
identify them. No counterexample to fullness under the \emph{entire}
working package is claimed. The designated small subobjects give a
colimit in $\K^{\mathrm{gen}}$ and in the ambient diagram category;
their colimit property for all cocones in $\mathcal C$ is not established.
Mixed-arrow factorization uses all the $H_\mu$, not just $H_{\mu_0}$.
\end{remark}

\begin{remark}[Type compatibility, formerly condition $(*)$]
\label{rem:working-orbit-invariance}
For $X,Y\in\mathcal C_\mu$, $Z\in\mathcal C_\nu$, $\mu\leq\nu$,
any $f\in H(X,Z)$, $g\in H(Y,Z)$, and bounded $\ba,\bb$,
\[
 \exists s:X\cong Y\ (s\ba=\bb)
 \quad\Longleftrightarrow\quad
 \exists u\in\operatorname{Aut}(Z)\ (uf\ba=g\bb).
\]
Indeed, arrow preservation identifies the two outer lifted types with
their respective images, and the orbit axiom applies at each level.
Thus there is no need to add $(*)$ separately. The lifting relation is
primary; its equivalence and orbit characterization remain substantive
axioms. In particular, an $H$-arrow agrees on each bounded tuple with
a system isomorphism at the same level. This is not automatic for
arbitrary structural or strong embeddings.
\end{remark}

\begin{remark}[Exhaustion is not an assertion about arbitrary partial maps]
Compatible diagram matches give partial maps $p_n:A_n\to Y$ with
$p_n\subseteq p_{n+1}$. Equality and atomic information make their union
a structural embedding once the $A_n$ exhaust $X$; fixed constants are
preserved when they occur. The partial domains need not be models.
The enumerated $B^v_n$ must occur \emph{within} the matched lists.
An unrelated resolution of each vertex is insufficient to invoke the
source criterion: a bounded set can meet all stages of a countable
exhaustion without being contained in any one stage.

The images $p[B^v_n]$ are nested strong submodels of $Y(v)$, and their
inclusions are strong by coherence. The chain and common-upper-bound
axioms of $\K$ show $p[X(v)]\leq_\K Y(v)$. This proves component
strongness, not a possible extra condition on $H$. At the intended
countable-cofinality levels, Lemma~\ref{lem:countable-exhaustion}
also resolves the earlier ordinal-indexing concern for nested lists.

In fact, a single bounded match from $X\in\mathcal C_\mu$ to
$Y\in\mathcal C_\nu$, $\mu\leq\nu$, extends to a whole $H$-arrow:
embed $X$ into some $Z\in\mathcal C_\nu$ by a presentation arrow,
use type invariance, and then use the orbit axiom to map $Z$ isomorphically
to $Y$ with the prescribed values. Thus the strong-image condition for
an enumerated $B^v_n$ follows from the full other axioms in this situation.
We retain it explicitly to match the source's criterion and its
verification obligations. The whole arrows extending different partial
matches need not agree outside their prescribed lists; their union is
not what is being taken.
\end{remark}

\begin{remark}[Warnings retained from the reading discussion]
\label{rem:working-template-arbitrary-orders}
\begin{enumerate}[label=\textup{(\roman*)},leftmargin=*]
\item Formal evaluation of $\Psi$ at an arbitrary linear order does not
establish membership in $\K^{\mathrm{gen}}$ or eligibility. Source
Lemma~5.4 concerns padded orders and uses additional arguments.
No skeleton or term-label expansion is silently named in the
structured-isomorphism relation.
\item System categoricity is stronger than vertex categoricity. For the
walking-arrow shape and pure sets, a bijective internal map and a proper
injection are nonisomorphic systems even with all vertices of size $\mu$.
The source verifies suitable system packages in Proposition~5.6 and
Theorems~8.6 and~9.3; Proposition~11.7 assumes the needed geometric-system
categoricity, and Proposition~5.8 takes categorical systems as input.
\item Rigid indexing well-orders do not prevent automorphisms of the
system reduct. Matching support patterns are compared through order
arrows into a common presentation, not automorphisms of the index order.
At one level the orbit axiom then gives a system automorphism.
\item These lifted types are not identified here with Galois types in
the entire generated AEC. No general amalgamation assumption for that
AEC is implicit. The comparison with placed diagrams after selected-level
amalgamation remains a separate reading question.
\end{enumerate}
\end{remark}
